% Options for packages loaded elsewhere
\PassOptionsToPackage{unicode}{hyperref}
\PassOptionsToPackage{hyphens}{url}
\documentclass[
  ignorenonframetext,
]{beamer}
\newif\ifbibliography
\usepackage{pgfpages}
\setbeamertemplate{caption}[numbered]
\setbeamertemplate{caption label separator}{: }
\setbeamercolor{caption name}{fg=normal text.fg}
\beamertemplatenavigationsymbolsempty
% remove section numbering
\setbeamertemplate{part page}{
  \centering
  \begin{beamercolorbox}[sep=16pt,center]{part title}
    \usebeamerfont{part title}\insertpart\par
  \end{beamercolorbox}
}
\setbeamertemplate{section page}{
  \centering
  \begin{beamercolorbox}[sep=12pt,center]{section title}
    \usebeamerfont{section title}\insertsection\par
  \end{beamercolorbox}
}
\setbeamertemplate{subsection page}{
  \centering
  \begin{beamercolorbox}[sep=8pt,center]{subsection title}
    \usebeamerfont{subsection title}\insertsubsection\par
  \end{beamercolorbox}
}
% Prevent slide breaks in the middle of a paragraph
\widowpenalties 1 10000
\raggedbottom
\AtBeginPart{
  \frame{\partpage}
}
\AtBeginSection{
  \ifbibliography
  \else
    \frame{\sectionpage}
  \fi
}
\AtBeginSubsection{
  \frame{\subsectionpage}
}
\usepackage{iftex}
\ifPDFTeX
  \usepackage[T1]{fontenc}
  \usepackage[utf8]{inputenc}
  \usepackage{textcomp} % provide euro and other symbols
\else % if luatex or xetex
  \usepackage{unicode-math} % this also loads fontspec
  \defaultfontfeatures{Scale=MatchLowercase}
  \defaultfontfeatures[\rmfamily]{Ligatures=TeX,Scale=1}
\fi
\usepackage{lmodern}
\ifPDFTeX\else
  % xetex/luatex font selection
\fi
% Use upquote if available, for straight quotes in verbatim environments
\IfFileExists{upquote.sty}{\usepackage{upquote}}{}
\IfFileExists{microtype.sty}{% use microtype if available
  \usepackage[]{microtype}
  \UseMicrotypeSet[protrusion]{basicmath} % disable protrusion for tt fonts
}{}
\makeatletter
\@ifundefined{KOMAClassName}{% if non-KOMA class
  \IfFileExists{parskip.sty}{%
    \usepackage{parskip}
  }{% else
    \setlength{\parindent}{0pt}
    \setlength{\parskip}{6pt plus 2pt minus 1pt}}
}{% if KOMA class
  \KOMAoptions{parskip=half}}
\makeatother
\setlength{\emergencystretch}{3em} % prevent overfull lines
\providecommand{\tightlist}{%
  \setlength{\itemsep}{0pt}\setlength{\parskip}{0pt}}
% Slide layout defaults for warning-clean scientific decks.
\usepackage{etoolbox}
\IfFileExists{xurl.sty}{\usepackage{xurl}}{}
\IfFileExists{seqsplit.sty}{\usepackage{seqsplit}}{\newcommand{\seqsplit}[1]{#1}}
\protected\def\breakseq#1{\seqsplit{#1}}
\protected\def\breaktt#1{\begingroup\ttfamily\seqsplit{#1}\endgroup}
\setlength{\emergencystretch}{6em}
\tolerance=5000
\hbadness=10000
\hfuzz=1pt
\setlength{\tabcolsep}{2pt}
\AtBeginEnvironment{longtable}{\tiny\renewcommand{\arraystretch}{0.86}\setlength{\tabcolsep}{1pt}}
\AtBeginEnvironment{tabular}{\tiny\renewcommand{\arraystretch}{0.86}\setlength{\tabcolsep}{1pt}}

% Natbib and cross-reference fallbacks — slides don't load natbib
% or cleveref, but combined-PDF manuscript prose may emit these
% commands. The fallback renders citations as a bracketed key list
% and cross-references as detokenized labels so slides don't fail on
% undefined control sequences or raw underscores. \providecommand is
% a no-op if packages load later.
\providecommand{\citep}[1]{[#1]}
\providecommand{\citet}[1]{#1}
\providecommand{\citealp}[1]{#1}
\providecommand{\citeauthor}[1]{#1}
\providecommand{\citeyear}[1]{#1}
\providecommand{\cref}[1]{\texttt{\detokenize{#1}}}
\providecommand{\Cref}[1]{\texttt{\detokenize{#1}}}

\newtheorem{proposition}{Proposition}
\newtheorem{hypothesis}{Hypothesis}
\usepackage{bookmark}
\IfFileExists{xurl.sty}{\usepackage{xurl}}{} % add URL line breaks if available
\urlstyle{same}
% fallback for those not using the hyperref driver hyperxmp:
\makeatletter
\@ifundefined{xmpquote}{\newcommand{\xmpquote}[1]{#1}}{}
\makeatother
\hypersetup{
  hidelinks,
  pdfcreator={LaTeX via pandoc}}

\author{\texorpdfstring{}{}}
\date{}

\begin{document}

\begin{frame}[allowframebreaks]{Field Overview: Disciplinary Structure
and Growth Dynamics \label{sec:field_overview}}
\protect\phantomsection\label{field-overview-disciplinary-structure-and-growth-dynamics}
Annual output in the Active Inference literature rose from 1 papers in
2005, reaching a peak of 123 papers in 2025---a transition from a niche
within theoretical neuroscience to a multi-disciplinary research program
spanning three primary domains and eight tracked categories. The corpus
start of 2005 was chosen to capture Energy-Based Model and variational
Bayesian antecedents \citep{dayan1995helmholtz, lecun2006tutorial} that
preceded the formal introduction of the Free Energy Principle in 2006
\citep{friston2006free} and its subsequent full elaboration
\citep{friston2010free}. Our corpus, extracted from arXiv, Semantic
Scholar, and OpenAlex and deduplicated to \(N = 817\) papers
(2005--2026), captures the breadth, tempo, and internal architecture of
this expansion (Figure \ref{fig:field_summary}).

\begin{figure}[htbp]
\centering
\includegraphics[width=0.9\textwidth]{../figures/field_summary.png}
\caption{Publication counts by domain ($N = 817$). Application domains (C1--C5) collectively account for the largest share of the corpus; Domain A2 (qualitative philosophy) is the largest single category, reflecting the FEP's broad theoretical reach.}
\label{fig:field_summary}
\end{figure}

\begin{block}{Corpus-Level Summary}
\protect\phantomsection\label{corpus-level-summary}
\begin{table}[htbp]
\centering
\caption{Corpus-level summary statistics for the Active Inference literature corpus ($N = 817$), spanning 2005--2026.}
\label{tab:corpus_summary}
\begin{tabular}{ll}
\toprule
\textbf{Metric} & \textbf{Value} \\
\midrule
Total papers & 817 \\
Year range & 2005--2026 \\
Peak year & 2025 \\
CAGR & 20.36\% \\
Active domains & 8 of 8 tracked (A1--A2, B, C1--C5) \\
\bottomrule
\end{tabular}
\end{table}

The CAGR of 20.36\% (measured as the annualised growth rate of yearly
publication volume between endpoint years 2005 and 2026) reflects
sustained field expansion; annual output accelerated markedly in the
late 2010s---from roughly 2018 onward (Figure \ref{fig:growth_curve}).
Sustained high output persisting into subsequent years suggests the
field has reached a mature production phase rather than experiencing a
transient spike. Citation network metrics are detailed in the dedicated
citation network analysis (see
\hyperref[sec:citation_network]{the citation network analysis}).

\begin{figure}[htbp]
\centering
\includegraphics[width=0.8\textwidth]{../figures/growth_curve.png}
\caption{Annual (bars) and cumulative (line) publication counts, 2005--2026 ($N = 817$, CAGR = 20.36\%). Annual output accelerates from the late 2010s (around 2018) to the 2025 peak. Moving average trendline (dashed), peak year, and median year annotated.}
\label{fig:growth_curve}
\end{figure}
\end{block}

\begin{block}{Domain Distribution}
\protect\phantomsection\label{domain-distribution}
Keyword-based classification assigns each paper to one of eight
categories across three domains:

\begin{table}[htbp]
\centering
\caption{Domain distribution across three tiers and eight categories ($N = 817$ papers). Classification uses hierarchical keyword matching with priority-based routing to minimize over-assignment to catch-all categories.}
\label{tab:domain_distribution}
\begin{tabular}{llcc}
\toprule
\textbf{Domain} & \textbf{Category} & \textbf{Papers} & \textbf{Percentage} \\
\midrule
A -- Core Theory & A1: Formal Theory & 64 & 7.8\% \\
 & A2: Qualitative Philosophy & 60 & 7.3\% \\
\midrule
B -- Tools & B: Tools \& Translation & 170 & 20.8\% \\
\midrule
C -- Applications & C1: Neuroscience & 161 & 19.7\% \\
 & C2: Robotics & 136 & 16.6\% \\
 & C3: Language & 58 & 7.1\% \\
 & C4: Psychiatry & 33 & 4.0\% \\
 & C5: Biology & 135 & 16.5\% \\
\bottomrule
\end{tabular}
\end{table}

The concentration of papers in A2 (qualitative philosophy and general
theory) reflects the broad scope of foundational FEP work (Figure
\ref{fig:subfield_distribution}). The priority-based classifier
mitigates over-assignment by routing papers with mathematical indicators
(theorems, proofs, equations, statistical formalism) to A1 before
falling back to A2, and by preferring specific application domains
(C1--C5) and tools (B) over both core-theory categories. Papers that
discuss FEP/AIF conceptually without mathematical formalism or
domain-specific vocabulary are correctly assigned to A2. This figure
should be read as a \emph{ceiling} on theoretical generality rather than
a literal measure of research focus---embedding-based classification
would likely redistribute some fraction into more specific categories.
That all eight categories are populated, including computational
psychiatry (C4) and formal theory (A1), indicates diversification beyond
the field's neuroscience origins.

\begin{figure}[htbp]
\centering
\includegraphics[width=0.8\textwidth]{../figures/subfield_distribution.png}
\caption{Domain distribution ($N = 817$). Classification uses hierarchical keyword matching against curated lists applied to titles and abstracts, capturing distinct methodological and domain-specific groupings.}
\label{fig:subfield_distribution}
\end{figure}

Detailed characterizations of each domain---including historical
context, growth trends, and open problems---are provided in the
supplementary domain analyses (see
\hyperref[sec:subfield_analyses]{the domain analyses}). Latent topic
structure, vocabulary analysis, and document embeddings are presented in
the text analytics section (see
\hyperref[sec:text_analytics]{the text analytics section}).
\end{block}

\begin{block}{Cross-Domain Comparison}
\protect\phantomsection\label{cross-domain-comparison}
\begin{table}[htbp]
\centering
\caption{Cross-domain comparison showing growth trajectories, maturity levels, key challenges, and representative publications for each of the eight tracked categories. Growth trends and maturity assessments are based on temporal publication patterns and evidence base depth.}
\label{tab:cross_domain_comparison}
\begin{tabular}{llcllll}
\toprule
\textbf{Domain} & \textbf{Category} & \textbf{Papers} & \textbf{Growth} & \textbf{Maturity} & \textbf{Key Challenge} & \textbf{Rep.\ Work} \\
\midrule
A & A1: Formal & 64 (7.8\%) & Growing & Mature & Math accessibility & \citep{sakthivadivel2023bayesian} \\
A & A2: Philosophy & 60 (7.3\%) & Stable & Mature & Catch-all absorption & \citep{friston2010free} \\
B & B: Tools & 170 (20.8\%) & Rapid & Growing & Deep RL benchmarks & \citep{fountas2020deep} \\
C & C1: Neuroscience & 161 (19.7\%) & Stable & Mature & Theory--neuroimaging gap & \citep{clark2013whatever} \\
C & C2: Robotics & 136 (16.6\%) & Growing & Growing & Embedded real-time & \citep{lanillos2021active} \\
C & C3: Language & 58 (7.1\%) & Emerging & Nascent & NLP model comparison & \citep{friston2020generative} \\
C & C4: Psychiatry & 33 (4.0\%) & Emerging & Nascent & Clinical translation & \citep{smith2021computational} \\
C & C5: Biology & 135 (16.5\%) & Rapid & Nascent & Empirical validation & \citep{kuchling2020morphogenesis} \\
\bottomrule
\end{tabular}
\end{table}

Three structural features emerge from the cross-domain comparison
(Figure \ref{fig:subfield_timeline}). First, no single legacy domain
dominates: Domain B (Tools \& Translation) accounts for 20.8\% of the
corpus, followed by C1 (Neuroscience) at 19.7\% and C2 (Robotics) at
16.6\%. Second, Domain A (Core Theory) aggregates 15.2\% collectively
(A1 + A2), while the emergent application frontiers (C3--C5) exhibit
accelerating growth. Third, A1's 64 papers understate its intellectual
influence---the mathematical formalisms developed in A1 shape
implementations across all domains.

\begin{figure}[htbp]
\centering
\includegraphics[width=0.9\textwidth]{../figures/subfield_timeline.png}
\caption{Stacked area chart of publications by domain, 2005--2026 ($N = 817$). A2 (qualitative philosophy) provides a large baseline; application domains C1--C5 show accelerating diversification from 2015 onward.}
\label{fig:subfield_timeline}
\end{figure}
\end{block}
\end{frame}

\end{document}
